\scanpage{105}
\mspage{52}
4. \uline{Sections \inserted{d'extensions,} et anneaux de valuation généraux.}

D'abord une révision de notations. Si $X$ est un schéma connexe, je \deleted{\unsure{dé|de}}note \deleted{$\mathcal E_X$}
\[
(1)\qquad \mathcal E_X=\Pi_1(X)
\]
son groupe fond. profini en tant que groupe extérieur, et si $\tilde X$ est un rev. universel profini de $\mathcal E_X$, par
\[
(2)\qquad \mathcal E_X^{(\tilde X)}=\Pi_1(X;\tilde X)=\operatorname{Aut}_X(\tilde X)
\]
son groupe fondamental précisé --- qui est un groupe profini. Si $X=\operatorname{Spec}(A)$, où $A$ est un anneau (le plus souvent un corps) je note $\mathcal E_A$, et $\mathcal E_A^{\tilde A}$ \inserted{$=\mathcal E_X^{\unsure{\tilde X|X}}$} si $\tilde A$ est une $A$-algèbre telle que $\operatorname{Spec}(\tilde A)=\tilde X$ soit un rev. universel de $X$ (ce qui le détermine à isom non unique près). Bien entendu, si $\unsure{\zeta|\xi}$ est un ``point géométrique'' de $X$, on note
\[
(3)\qquad \mathcal E_X^{\unsure{\zeta|\xi}}=\Pi_1(X,\unsure{\zeta|\xi})=\mathcal E_X^{\tilde X(\unsure{\zeta|\xi})},
\]
où $\tilde X(\unsure{\zeta|\xi})$ est le \unsure{rev|rev.} universel de $X$ pointé en $\unsure{\zeta|\xi}$. Le choix
\margin{104}
